% Volume V: Admissibility and Continuation
% Part X. Public Institutions as Epistemic Ecologies
% Chapter 53: Separation of Powers
% Spine: proof-spines/ch053-spine.md

\chapter{Separation of Powers}
\label{ch:ch053}

Separation of powers is treated here as structured non-identical observation. Different institutions possess different information, incentives, temporal horizons, and error modes, so they do not process appearances in the same way even when they address the same public matter.

The goal is not perfect independence, but prevention of correlated collapse. Institutional redundancy succeeds when disagreement remains possible long enough to expose hidden assumptions, local blind spots, or premature closures that a single integrated line of authority would normalize. The separation of powers is therefore an epistemic ecology before it is only a constitutional slogan.
